% =====================================================================
% Appendix B --- The model, its closures and the formal results
% ---------------------------------------------------------------------
% STRUCTURE (2026-10-22): notation folded into B.1 "Production and
% consumption"; the "Purpose and summary" subsection removed; the shared
% blocks renamed B.4 "The closure systems"; all subsection labels are app:*
% (was: sec:reduction, sec:equiv, sec:sectoral, sec:gauge, sec:scope), and
% the section carries ONE label, app:proofs (the aliases app:model and
% app:closure-focs were retired; Section 4 now points at app:proofs).
% Admissibility paragraph removed; B.10 compressed to the evidentiary minimum.
% MERGED 2026-10-18 from three appendices that overlapped (closure FOCs, model
% formulation, proofs). One equation label set only; proof-specific equations
% (eq:int, eq:cons, eq:leontiefblocks, eq:affine, eq:bvec, eq:finstrument)
% remain where first used; equations run in one series.
% REGISTER: no ADR numbers, run ids, generation tags, probe names, file or
% registry paths; measured values are stated as values. Keep it that way.
% =====================================================================

\section{The model, its closures and the formal results}\label{app:proofs}
%NOTE (recorded 2026-10-23 at the author's request): the author doubts that what
%remains of this appendix carries its weight for a Metroeconomica audience and
%considers Appendix A sufficient on its own. What is kept here is either exact
%(Lemmas lem:price and lem:affine; Proposition prop:nest) or
%measured (Table~\ref{tab:sectoral} with the sectoral readings). Cut in this round: the
%classification corollary, the off-slice remark and the interpolation remark ---
%the last because the limit it asserted does not hold. Candidate for deletion
%before resubmission if no referee asks for proofs.

\noindent
This appendix collects the formal apparatus behind the closure matrix: the
notation and the functional forms that solve the producer and household
problems of Section~\ref{subsec:optimization}
(Section~\ref{app:model-kernel}); the five labor-market closures and the
three financing closures (Sections~\ref{app:labour-closures}
and~\ref{app:financing-closures}); and the formal results --- the closure
systems, the demand-only reduction and the rigidity result
(Proposition~\ref{prop:nest};
Sections~\ref{app:closure-systems}--\ref{app:sectoral}). The
intermediate-bill and external-account accounting is in
Appendix~\ref{app:abill}.


\subsection{Production and consumption}\label{app:model-kernel}

\paragraph{Notation.} $L_i$ is employment in sector $i$ and $L=\sum_i L_i$ total employment; a bar marks the baseline, so $\bar L_i$ is baseline sectoral and $\Lbar$ baseline total employment. $w$ is the common wage of the mobile closures, $w_i$ the
sector-specific wage where labor markets are sectoral, and $\bar w$ the baseline real-wage anchor. Real wages are computed by deflating with the consumption price index $\mathcal{P}(p)$, the CES price of the household's bundle; a subscript $0$ marks
baseline values. $L_i^{\mathrm{cm}}(p,y_i,w_i)$ is the labor demand that
solves the firm's cost-minimization problem
of Section~\ref{subsec:optimization} at prices $p$ and output $y_i$. The
allocation parameter $\eta$ takes only its endpoints: $\eta=1$ lets labor
move to its cost-minimizing use, $\eta=0$ freezes the baseline allocation of labor.
The labor-supply elasticity $\eta_s$ governs how strongly sectoral employment
responds to the real wage; in principle, it may differ across sectors ($\eta_{s,i}$), but the paper uses a
common value $\eta_s$ throughout. Finally, $\theta$, $\epsilon$ and $\sigma$ are the
substitution elasticities of Section~\ref{subsec:optimization} (among
intermediates, between labor and intermediates, and in consumption); their
limit $0^{+}$ is fixed proportions. Table~\ref{tab:parameters} collects these elasticities and the two parameters related to the extensions in Section~\ref{sec:beyond-demand-only} for convenience.

% =====================================================================
% Table --- the elasticities of the comparison (Sec. 4.2)
% ---------------------------------------------------------------------
% Requested by reviewer 1 (R1.3: "including elasticity values in a table").
% Values are read from the executed design experiments/designs/
% matrix_5x3_v11.toml and from the manifests of the reported cells
% (runs/matrix_5x3-v11-*): sigma, theta, epsilon, eta are the design values;
% eta_s is the bracket printed in Table 4; the two extension parameters are
% the ones the reported tables of Section 6 use. None of these is estimated
% from the input--output table.
% =====================================================================
\begin{table}[t]
\centering
\small
\begin{tabular}{@{}l p{5.9cm} l l@{}}
\toprule
Parameter & Meaning in the kernel & Value & Set in \\
\midrule
$\sigma$ & substitution across consumption goods & $0.9$ & all cells \\
$\theta$ & substitution across intermediates & $0.5$ & all cells \\
$\varepsilon$ & labor--intermediate nesting & $0.5$ & all cells \\
$\eta$ & intersectoral labor-allocation selector & $\{0,1\}$ & labor closure \\
$\eta_s$ & sectoral labor-supply elasticity with respect to the real wage & $\{0.5,1,2\}$ & elastic-supply column \\
\addlinespace

$\delta$ & capacity door: utilization elasticity of productivity & $\{0,0.05\}$ & Section~\ref{subsec:pricing-capacity} \\
$\kappa$ & markup door: margin--utilization slope & $\{0.25,0.5\}$ & Section~\ref{subsec:pricing-markup} \\

\bottomrule
\end{tabular}
\par\smallskip
\parbox{\textwidth}{\footnotesize
\textit{Note:} $\eta=0$ fixes employment in each sector at its baseline level,
with sector-specific wages; $\eta=1$ permits cost-minimizing labor
reallocation at a common wage.}
\caption{Parameter values for the baseline comparison and pricing extensions.}
\label{tab:parameters}
\end{table}


On the financing side, $c_i^h$ is the household purchase of good $i$ and
$E^h$ its consumption expenditure, so $\sum_i p_i c_i^h=E^h$ is the
household budget of equation \eqref{eq:utility-max}. $W=\sum_i w_i L_i$ is the wage bill, $T_G(p)=\sum_i p_i g_i^G$ the baseline tax cost of government purchases, and $g_i$ the program bundle (real public investment in sector $i$) with current-price cost
$\sum_i p_i g_i$. F1 changes the composition of $E^h$, while under F2 the program's cost is financed by the wage-income tax at the endogenous rate $\tau$ defined in~\eqref{eq:cons}; under F3 it is booked as external program financing $B_{\mathrm{gov}}$. The net external transfer $F$ closes the external account; it is endogenous in the mobile regime and identically zero in the wage-pinned regime.

Each sector's technology is a nested CES with constant returns: labor and an intermediate composite substitute with elasticity $\epsilon$, and intermediate inputs substitute with elasticity $\theta$ inside the composite. Cost minimization, equation \eqref{eq:cost-min}, gives the intermediate price index and the unit cost
\begin{align}
P_i(p) &= \Big(\textstyle\sum_{u=1}^N \Omega_{iu}\,p_u^{1-\theta}\Big)^{1/(1-\theta)},
\label{eq:model-ip}\\
c_i(p,w_i;A_i) &= \Big(A_i^{\epsilon-1}\big(\alpha_i\,w_i^{1-\epsilon}
  + (1-\alpha_i)\,P_i^{1-\epsilon}\big)\Big)^{1/(1-\epsilon)},
\label{eq:model-cost}
\end{align}
with the Cobb--Douglas limits at $\theta=1$ and $\epsilon=1$. Because
$\sum_u\Omega_{iu}=1$, the unit cost equals one at the calibration point
$(p,w,A)=(\one,\one,\one)$. Zero profit requires
\begin{equation}
p_i = c_i(p,w_i;A_i) \qquad \forall\, i.
\label{eq:zp}
\end{equation}
Shephard's lemma, evaluated at zero profit ($p_i = c_i$), gives the
cost-minimizing labor demand
\begin{equation}
L_i^{\mathrm{cm}}(p,y_i,w_i)
 = \alpha_i A_i^{\epsilon-1} y_i
   \left(\frac{p_i}{w_i}\right)^{\epsilon},
\label{eq:closure-labour-demand}
\end{equation}
where $y_i$ is gross output, $A_i$ productivity and $\alpha_i$ the calibrated
labor share. Demand for intermediate goods follows in the same way from the
composite~\eqref{eq:model-ip} and is charged with the domestic bill
(see \eqref{eq:int}).

The household has CES preferences with elasticity $\sigma$ and calibrated
weights $\beta_i$. The consumption price index, which is also the
numeraire, is
\begin{equation}
  \mathcal{P}(p) = \left(\sum_{i=1}^N \beta_i p_i^{1-\sigma}\right)^{1/(1-\sigma)},
  \qquad \sum_{i=1}^N \beta_i = 1,
  \label{eq:labour-cpi}
\end{equation}
with the Cobb--Douglas limit at $\sigma = 1$. Given household consumption
expenditure $E^h$, utility maximization gives
\begin{equation}
c_i^h = E^h\frac{\beta_i p_i^{-\sigma}}{\sum_j\beta_j p_j^{1-\sigma}}.
\label{eq:model-household}
\end{equation}
Goods-market clearing for each of the $N$ sectors reads
\begin{equation}
y_i = \mathit{int}_i + (1-m_i)\big(c_i^h + g_i^G + \mathit{inv}_i + g_i\big) + X_i,
\label{eq:model-clearing}
\end{equation}
where $m_i$ is the sectoral import margin, $g^G$ baseline government
purchases, $\mathit{inv}$ non-program investment, $X$ exports and $g$ the
program bundle (zero under F1). Every closure below enforces all $N$
clearings~\eqref{eq:model-clearing}.

\subsection{The five labor closures}\label{app:labour-closures}

Let $L_i$ denote employment in sector $i$, $\bar{L}_i$ its baseline value,
and $\Lbar = \sum_{i=1}^N \bar{L}_i$ total baseline employment. Sectoral
wages are $w_i$; with mobile labor they equal one common wage $w$. All
closures use the same labor demand~\eqref{eq:closure-labour-demand}; they
differ in the wage and labor-supply restrictions imposed on it.

\paragraph{The no-reallocation closure: immobile labor.}
The no-reallocation row fixes employment in every sector and solves for a
sector-specific wage to clear each labor market:
\begin{equation}
L_i = L_i^{\mathrm{cm}}(p,y_i,w_i) = \bar{L}_i,
\qquad i=1,\ldots,N, \qquad \eta = 0.
\label{eq:closure-bf}
\end{equation}
Thus $\sum_i L_i = \Lbar$, but wages may differ across sectors.
The allocation parameter $\eta$ selects only the endpoints $\{0,1\}$;
its mobile endpoint $\eta = 1$ is the standard CGE closure below, not an
additional closure.

\paragraph{The standard CGE closure: mobile labor at full employment.}
Labor reallocates across sectors at one flexible wage, while total
employment remains fixed:
\begin{equation}
L_i = L_i^{\mathrm{cm}}(p,y_i,w),
\qquad \sum_{i=1}^N L_i = \Lbar,
\qquad \eta = 1.
\label{eq:closure-alpha}
\end{equation}

\paragraph{The elastic labor-supply closure: sectoral supply curves.}
We use the uniform-elasticity sectoral form: each sector has its own wage
and supply curve, with one common labor-supply elasticity $\eta_s$:
\begin{equation}
L_i = L_i^{\mathrm{cm}}(p,y_i,w_i)
 = \bar{L}_i \left[\frac{w_i/\mathcal{P}(p)}{w_{i0}/\mathcal{P}_0}\right]^{\eta_s},
\qquad \eta_s \geq 0, \qquad i=1,\ldots,N.
\label{eq:closure-beta}
\end{equation}
Equal supply elasticities do not impose equal wages: $N$ separate labor
markets determine $w_i$. As explained in Appendix~\ref{app:equivalences}, $\eta_s$ is a labor-supply
elasticity and not the allocation parameter $\eta$; yet setting $\eta_s = 0$
reproduces the no-reallocation row exactly. The calibrated real-wage anchors are $w_{i0}/\mathcal{P}_0 = 1$. The deflator remains the common consumption price index, not the sector's own goods price $p_i$. The supply curves are reduced forms without a
labor--leisure income effect. The headline panels use $\eta_s = 0.5$;
the elasticity ladder is reported in Table~\ref{tab:sectoral}. 

\paragraph{The employment-adjustment closure: wage pinned, employment uncapped.}
No labor-market equation is imposed, so employment is demand-determined and the
money wage serves as the numeraire:
\begin{equation}
w = 1 \quad(\text{numeraire}),
\qquad F = 0,
\qquad L_i = L_i^{\mathrm{cm}}(p,y_i,w),
\qquad \eta = 1 .
\label{eq:closure-gamma}
\end{equation}
Unlike the standard CGE closure, this closure does not impose
$\sum_i L_i = \Lbar$: employment absorbs whatever labor equilibrium
production requires; no employment constraint is imposed at all. The real
wage is an outcome and not an instrument --- at every solution
$w/\mathcal{P}(p) = 1/\mathcal{P}(\pi(A))$, a function of technology alone
(Corollary~\ref{cor:realwage}) --- so the pin $w=1$ is a numeraire rather than a
wage-setting rule. 

\paragraph{The Leontief closure: the fixed-coefficient corner.}
The Leontief closure keeps the employment rule of the preceding closure ---
pinned money wage, uncapped and demand-determined employment, $F = 0$ ---
while taking the substitution elasticities to their Leontief limits:
\begin{equation}
w = 1, \qquad F = 0, \qquad \eta = 1,
\qquad (\theta,\epsilon,\sigma) \longrightarrow (0^+,0^+,0^+).
\label{eq:closure-delta}
\end{equation}
The Leontief closure is therefore the fixed-coefficient corner of the
employment-adjustment closure rather than an independent
labor-market mechanism. The two rows share their solution only while prices
sit at the baseline, that is, at the calibration technology under a
demand-only impulse (Lemma~\ref{lem:affine}); as soon as prices move,
they separate (Section~\ref{sec:beyond-demand-only}).
% NOTE (2026-10-23): reworded from "The corner is exact at the calibration
% technology". That phrase was broader than the demand-only proposition assumes: the
% markup runs of Section sec:beyond-demand-only sit at unchanged technology
% (A = 1) and still separate the two rows, since demand enters the price
% block there. The exact boundary is prices at the baseline.

\subsection{The three financing closures}\label{app:financing-closures}
% Programme-scale F1 tilt: experiments/run.jl, f1_tilt_weights.

Each labor closure is crossed with three ways of accommodating the same
program scale and sectoral incidence. Let $g_i = G_0\psi_i$ denote the
reference program bundle, where $G_0$ is its baseline-price value and
$\psi_i \geq 0$ with $\sum_i\psi_i = 1$ describes its sectoral incidence;
a $k$-fold program is $kg$. F2 and F3 add this
bundle to final demand; F1 instead uses its scale and incidence to redirect
household spending.

Baseline government purchases $g_i^G$ remain tax-financed in all three
cases. Define the wage bill and their current-price tax cost as
\begin{equation}
W = \sum_{i=1}^N w_i L_i,
\qquad T_G(p) = \sum_{i=1}^N p_i g_i^G.
\label{eq:financing-baseline-budget}
\end{equation}
With a common wage, $W = w\sum_i L_i$. Let $E$ denote after-tax household
resources, including the net external transfer $F$ from the rest of the
world,\footnote{$F$ is the variable that clears the external account; the
labels F1--F3 name financing closures, not a variable.} and let $s$ be the
calibrated saving rate. Household consumption expenditure is
$E^h = (1-s)E$. Gross household purchases $c_i^h$ exhaust this budget:
\begin{equation}
\sum_{i=1}^N p_i c_i^h = E^h,
\qquad c_i^{\mathrm{dom}} = (1-m_i)c_i^h.
\label{eq:financing-household-budget}
\end{equation}
Only the domestic content enters domestic goods-market
clearing~\eqref{eq:model-clearing}; likewise, the program contributes
$(1-m_i)g_i$ under F2 and F3.

\paragraph{F1: budget-neutral preference reallocation.}
F1 changes the household's CES weights without adding a public-demand
bundle. For positive preference shifters $d_i$, the weights and gross
household demand become
\begin{align}
\widetilde{\beta}_i
 &= \frac{\beta_i d_i}{\sum_j\beta_j d_j},
\qquad \sum_i\widetilde{\beta}_i = 1,
\label{eq:closure-f1-weights}\\
c_i^h
 &= E^h\frac{\widetilde{\beta}_i p_i^{-\sigma}}
 {\sum_j\widetilde{\beta}_j p_j^{1-\sigma}},
\qquad E = W-T_G(p)+F.
\label{eq:closure-f1}
\end{align}
The program-scale tilt is $d_i = 1+kG_0\psi_i^+/c_{i0}^h$ in sectors
with positive baseline household purchases, and $d_i=1$ otherwise;
$\psi^+$ is the program incidence restricted to those sectors and
renormalized to sum to one. The CES normalizer preserves
\eqref{eq:financing-household-budget} at every price vector. Thus increased
spending shares in program sectors are paid for by reduced shares
elsewhere, conditional on the household budget; equilibrium income need
not remain unchanged. Hence, strictly speaking, F1 is a demand-repurposing experiment emerging from either preference shifts or imposed regulation, not an autonomous investment injection or multiplier.

\paragraph{F2: tax-financed public investment.}
F2 adds the real program bundle $g$ and levies an additional household
tax that covers the program at equilibrium prices:
\begin{equation}
T(p) = \sum_{i=1}^N p_i g_i,
\qquad E = W-T_G(p)-T(p)+F,
\qquad B_{\mathrm{gov}}=0.
\label{eq:closure-f2}
\end{equation}
Total domestic tax revenue $T_G(p)+T(p)$ therefore balances the purchases
$g^G+g$, priced at current rather than baseline prices. The tax is
implemented proportionally as $\tau=[T_G(p)+T(p)]/W$ and the tax base moves with the wage bill $W$. In essence, the program has an explicit domestic budget counterparty and reduces household resources through $\tau$.

\paragraph{F3: externally financed public investment.}
F3 adds the same real bundle but imposes no additional domestic program
tax. Its current-price cost is booked as external program financing:
\begin{equation}
B_{\mathrm{gov}} = \sum_{i=1}^N p_i g_i,
\qquad E = W-T_G(p)+F.
\label{eq:closure-f3}
\end{equation}
The booking $B_{\mathrm{gov}}$ is distinct from the net external transfer
$F$ entering household resources after tax, although their economic interpretation is similar, namely some form of external debt. $F$ is solved endogenously in the no-reallocation, the standard CGE and both elastic-supply forms, but is fixed at zero in the two wage-pinned closures, the employment-adjustment closure and its Leontief corner. $B_{\mathrm{gov}}$, by contrast, is set endogenously to match the size of the
spending program.
\begin{equation}
S + T_{\mathrm{int}} + M - (I+X) = F+B_{\mathrm{gov}},
\qquad S=sE,
\label{eq:financing-external-account}
\end{equation}
where $I$ is non-program investment, $X$ exports, $M$ total imports
(including intermediate and program imports), and $T_{\mathrm{int}}$
the separately booked product-tax leak on intermediate use; all are
current-price values. The net external position is therefore $F+B_{\mathrm{gov}}$, not the program
cost alone. F3 is an accounting open-economy closure, not a monetary mechanism: the static model cannot distinguish external borrowing from a transfer and does not model interest payments, repayment or exchange-rate adjustment.

\subsection{The closure systems}\label{app:closure-systems}

We keep the notation above and record the data-calibrated objects the
systems below use: domestic bill coefficients $a_i$, the value of domestic
intermediates used per unit of gross output of sector $i$, the input-share matrix $\Omega$ ($\Omega_{iu}$ the share of input $u$ in
sector $i$'s bundle, rows summing to one), sector import margins $m_i$, the
saving rate $s$, baseline government demand $g^G$, non-program investment
$inv$, exports $X$, and the labor bar $\Lbar = \sum_i \bar L_i = 1$ in
baseline income units (GDP $=1$).

\paragraph{Price block.} Prices and technique are those of the formulation
above: the intermediate price index is \eqref{eq:model-ip}, the unit cost
\eqref{eq:model-cost}, with the Cobb--Douglas limit
$c_i = w^{\alpha_i} P_i^{\,1-\alpha_i}/A_i$ at $\epsilon = 1$, and zero
profit imposes \eqref{eq:zp} for every sector. At the calibration point $(p,w,A) = (\one,1,\one)$ the unit cost equals one
exactly, because $\sum_u \Omega_{iu} = 1$ makes $P_i = 1$; this is the
calibration identity used throughout.

\paragraph{Quantity block.} With full cost-minimizing allocation
($\eta = 1$) the marginal-product condition inverts to the labor demand
\eqref{eq:closure-labour-demand}. At $\eta = 0$ the allocation is the
frozen baseline, $L_i = \bar{L}_i$; under the sectoral-wage closure the wage
$w$ in \eqref{eq:closure-labour-demand} is then a per-sector $w_i$, solved from
$L^{cm}_i(p, y_i, w_i) = \bar{L}_i$ jointly with \eqref{eq:zp}.
Intermediate demand is charged with the \emph{domestic}
bill:
\begin{equation}
int_i = p_i^{-\theta} \sum_j \Omega_{ji}\,
        p_j^{\epsilon} A_j^{\epsilon-1} P_j^{\theta-\epsilon}\, a_j\, y_j.
\label{eq:int}
\end{equation}
Household demand is CES with CPI deflation and the financing weights
$ds$ ($ds = 1$ except under F1, where $ds = d$, the preference tilt):
\begin{equation}
c_i = (1-s)(1-m_i)\,\beta_i\, ds_i\; E\; p_i^{-\sigma} \Big/
      \textstyle\sum_j \beta_j ds_j\, p_j^{1-\sigma},
\qquad
E = (1-\tau)\, w \textstyle\sum_i L_i + F,
\label{eq:cons}
\end{equation}
with the wage-income tax rate $\tau = p\!\cdot\! g^G / (w \sum_i L_i)$ under
F1/F3 and $\tau = p\!\cdot\!(g^G + g)/(w \sum_i L_i)$ under F2, and $F$ the
net external transfer entering expenditure after tax.
At $\eta = 0$ the wage bill is
$\sum_i w_i \bar{L}_i$ in \eqref{eq:cons}. Saving
$sE$ leaks; the import content is charged margin by margin. Goods-market
clearing is \eqref{eq:model-clearing}, with $g_i$ present only under the
additive financings F2/F3 ($g_i = 0$ under F1).

\paragraph{The four regimes.} The mobile systems solve for
$X = (p, y, w, F) \in \mathbb{R}^{2N+2}$ with the consumption price index
\eqref{eq:labour-cpi} as numeraire, and close with the labor equation
\eqref{eq:closure-alpha} in the standard CGE closure and
\eqref{eq:closure-beta} in the elastic-supply closure (where
$w_0/\mathcal{P}_0 = 1$). The wage-pinned systems solve for $X = (p, y) \in \mathbb{R}^{2N}$ with the
wage pinned $w \equiv 1$ and $F \equiv 0$; employment is a post-solve
outcome. The employment-adjustment closure runs them at the baseline elasticities,
the Leontief closure at the Leontief
corner $(\theta_\delta,\epsilon_\delta,\sigma_\delta) = (10^{-4},10^{-4},
10^{-4})$. The no-reallocation closure is the sectoral-wage system at
$\eta = 0$: the
allocation stays frozen at the baseline, $L_i = \bar{L}_i$, and $N$ sectoral
wages $w_i$ are solved from each sector's own first-order condition
$L^{cm}_i(p, y_i, w_i) = \bar{L}_i$ jointly with zero profit, with unknowns
$[\,p;\; y;\; w(1{:}N);\; F\,]$ and the external transfer $F$ free. All rows
in the matrix use $\eta = 1$ except the no-reallocation closure.

\paragraph{What the transfer does.} In the mobile regime the labor equation
selects the transfer. Writing $y(F) = (I-G)^{-1}(b_0 + f\,F)$ with
$f = (1-s)\,\mathrm{diag}(1-m)\,\tilde\beta(ds)$, the pin
$\alpha^{\!\top}y = \Lbar$ gives
\begin{equation}
F^{\ast} = \frac{\Lbar - \alpha^{\!\top}(I-G)^{-1} b_0}
               {\alpha^{\!\top}(I-G)^{-1} f}.
\label{eq:finstrument}
\end{equation}
This instrument is why the mobile F2 and F3 columns are real-neutral: F2
charges the program through the wage-income tax \eqref{eq:closure-f2} ($E = \alpha^{\!\top}y -
\sum g^G - \sum g + F$), F3 through the external account \eqref{eq:closure-f3} ($E = \alpha^{\!\top}y
- \sum g^G + F$), and the two $b$-vectors differ by exactly
$f\,\sum g$. The pin \eqref{eq:finstrument} therefore yields
$F^{\ast}_{F3} = F^{\ast}_{F2} - \sum g$, identical real allocations and an
identical net external position $F + B_{\mathrm{gov}}$ (measured:
$F_{F3} - F_{F2} = -G_0$ in every mobile row). Where no labor equation remains,
nothing determines $F$, which is why the wage-pinned rows fix it at zero.

\subsection{The demand-only reduction}\label{app:reduction}

\begin{lemma}[Price block demand-invariance]\label{lem:price}
The zero-profit system \eqref{eq:zp} is homogeneous of degree one in
$(p,w)$ jointly and contains no demand variable. Consequently, for any
technology $A$, every solution satisfies $p = w\,\pi(A)$, where
$\pi(A) \in \mathbb{R}^N_{++}$ solves $c_i(\pi,1;A) = \pi_i$ for all $i$;
relative prices and the ratio $p/w$ are functions of $A$ alone. Under the
CPI numeraire, $w = 1/\mathcal{P}(\pi(A))$ and $p = \pi(A)/\mathcal{P}(\pi(A))$; under the
wage pin, $w = 1$ and $p = \pi(A)$. At $A \equiv \one$ the calibration
identity gives $\pi = \one$, hence in both gauges
\begin{equation}
p = \one, \qquad w = 1, \qquad \mathcal{P} = 1,
\end{equation}
independently of the demand vector, of the financing closure, and of the
shock size. The lemma is a statement about the
\emph{single-wage} price block: with a per-sector wage vector (the
$\eta = 0$ endpoint) homogeneity in $(p, w)$ no longer holds
sector by sector and the argument does not apply.
The statement is scoped to the branch the numeraire
selects: $\pi(A)$ denotes the normalized solution the kernel's solves return,
and the equivalences below compare the systems on that branch (the CES price
system at $\theta<1$ with self-loops admits spiral branches in principle; no
executed cell visits one).
\end{lemma}

\begin{proof}
Homogeneity: $c_i(tp, tw; A) = t\, c_i(p,w;A)$ for $t>0$, by inspection of
\eqref{eq:model-cost}--\eqref{eq:model-ip}. Write $q = p/w$. Then \eqref{eq:zp} is
equivalent to $q_i = c_i(q,1;A)$ for all $i$, i.e.\ $q = \pi(A)$, and
the solution branch of \eqref{eq:zp} selected by the
numeraire is the ray $\{(w\,\pi(A), w)\}$. The
CPI \eqref{eq:labour-cpi} is homogeneous of degree one in $p$, so
$\mathcal{P}(w\,\pi) = w\,\mathcal{P}(\pi)$; the numeraire $\mathcal{P} = 1$ or the pin $w = 1$
selects the scale. At $A \equiv \one$, $\pi = \one$ solves
\eqref{eq:zp} at $w = 1$ because $c_i(\one,1;\one) = 1$ (the calibration
identity), and $\mathcal{P}(\one) = 1$ because $\sum_i \beta_i = 1$.
\end{proof}

\begin{corollary}[Real-wage invariance]\label{cor:realwage}
At every solution of the shared block,
\begin{equation}
\frac{w}{\mathcal{P}(p)} = \frac{1}{\mathcal{P}(\pi(A))},
\end{equation}
a function of technology alone. At $A \equiv \one$ this equals $1 =
w_0/\mathcal{P}_0$: the equilibrium real wage sits at its anchor for every demand
vector, every financing closure and every shock magnitude.
\end{corollary}

%TODO: Ich glaub das kann weg.
%\begin{remark}[What the wage-pinned system implements]\label{rem:empadj}
%Dropping the labor-supply equation and fixing $w = 1$ leaves employment
%demand-determined at the real wage the price block delivers:
%$L_i = L_i^{\mathrm{cm}}(p, y_i, w)$ with no aggregate constraint and no
%wage-setting rule. We therefore call it an \emph{employment-adjustment}
%closure, wage-pinned: employment is the absorbing margin, the net external
%transfer is fixed at $F = 0$, and the real wage is an outcome that coincides
%with its baseline level only while relative prices stay at the calibration point. The label separates the
%closure's two substantive components --- the fixed transfer and the uncapped
%employment margin --- from its unit convention $w = 1$, which carries no
%economics (Corollary~\ref{cor:nogo}).
%\end{remark}

\begin{lemma}[Elasticity-free affine quantity block]\label{lem:affine}
Fix $\eta = 1$ and $A \equiv \one$, and evaluate the quantity block at the
demand-only prices $(p,w) = (\one,1)$. Then for \emph{every} value of
$(\theta,\epsilon,\sigma)$:
\begin{align}
int &= \Omega^{\!\top}\mathrm{diag}(a)\, y, &
L &= \mathrm{diag}(\alpha)\, y, &
\sum_i L_i &= \alpha^{\!\top} y,
\label{eq:leontiefblocks}\\
c &= (1-s)\,\mathrm{diag}(1-m)\,\tilde\beta(ds)\; E, &
E &= \alpha^{\!\top} y - \textstyle\sum g^G + F
   \;\; (\text{F1/F3}), \nonumber
\end{align}
where $\tilde\beta(ds) = (\beta \odot ds)\,/\, \one^{\!\top}(\beta \odot
ds)$ and $E = \alpha^{\!\top} y - \sum g^G - \sum g + F$ under F2. Hence the
clearing system \eqref{eq:model-clearing} is exactly affine,
\begin{equation}
y = G\, y + b(F),
\qquad
G = \Omega^{\!\top}\mathrm{diag}(a)
  + (1-s)\,\mathrm{diag}(1-m)\,\tilde\beta(ds)\,\alpha^{\!\top},
\label{eq:affine}
\end{equation}
\begin{equation}
b(F) = (1-m)\odot\big( g^G + inv + g \big) + X
     + (1-s)\,\mathrm{diag}(1-m)\,\tilde\beta(ds)
       \big( F - \textstyle\sum g^G - \mathbf{1}_{F2}\textstyle\sum g \big),
\label{eq:bvec}
\end{equation}
with $g$ present only under the additive financings
F2/F3 and $\mathbf{1}_{F2}$ selecting the tax-financed column (the bracket is
$F-\sum g^G$ under F1/F3 and $F-\sum g^G-\sum g$ under F2).
No elasticity parameter enters $G$ or $b$. At
$\eta = 0$ the aggregate is still pinned, $\sum_i L_i = \Lbar$, but the
unit-price reduction below describes only the \emph{retired} common-wage
pin: there $E$ did not depend on $y$ and the household round dropped out,
$G_{\mathrm{nr}} = \Omega^{\!\top}\mathrm{diag}(a)$ (the no-reallocation
row). Under the sectoral-wage
closure the wages $w_i$ move with the demand composition, the prices leave
the unit point, and this affine form does not describe the $\eta = 0$ row.
\end{lemma}

\begin{proof}
At $p = w = A = \one$: $P_j = 1$ in \eqref{eq:int}, so all elasticity
factors $p^{\epsilon}$, $P^{\theta-\epsilon}$, $p^{-\theta}$ equal $1$ and
\eqref{eq:int} reduces to the fixed-coefficient round
$\sum_j \Omega_{ji} a_j y_j$; \eqref{eq:closure-labour-demand} reduces to
$L_i = \alpha_i y_i$ since every exponent collapses at unit prices; in
\eqref{eq:cons} the terms $p_i^{-\sigma}$ and the normalizer reduce to
$1$ and $\sum_j \beta_j ds_j$, giving the fixed spending weights
$\tilde\beta(ds)$; and the tax rate $\tau = \sum g^G / \alpha^{\!\top}y$
turns $E$ into the affine form shown. Collecting terms in
\eqref{eq:model-clearing} gives \eqref{eq:affine}--\eqref{eq:bvec}: the
intermediate round contributes the columns
$\sum_i \Omega_{ji} a_j = a_j$ (row sums of $\Omega$ equal $1$) and the
household round contributes $(1-s)\sum_i (1-m_i)\tilde\beta_i(ds)\,
\alpha_j$ per unit of $y_j$.
\end{proof}

\begin{remark}[Nonsubstitution]\label{rem:samuelson}
Lemma~\ref{lem:price} and Lemma~\ref{lem:affine} are the production side of
the Samuelson nonsubstitution theorem: with one primary factor, constant
returns and no joint production, the cost-minimizing technique --- and hence
relative prices --- is independent of the demand vector, and final demand
allocates quantities only (Samuelson 1951). The kernel's nested-CES
structure is a smooth parametrization of that corner; the demand-only matrix
therefore inherits the theorem's rigidity: this is the model-class
property behind the coincidences. The IO endpoint makes the rigidity
literal --- the Robinson (2006) fixed-price multiplier --- but, as
Lemma~\ref{lem:affine} shows, it is not an approximation:
the same candidate equilibrium $(p,y)=(\one,y^{\ast})$
solves the wage-pinned system at every elasticity triple in the wage-pinned
gauge.
\end{remark}

\subsection{The elastic labor supply model and its two ends}\label{app:sectoral}

The equivalences stated in Appendix~\ref{app:equivalences} are coincidences
\emph{within} the executed matrix. The sectoral elastic labor supply model extends beyond it, and
both of its ends are pinned: at $\eta_s=0$ it coincides exactly with the
matrix's no-reallocation row (Proposition~\ref{prop:nest}), and as
$\eta_s\to\infty$ it has a limit that aligns with the employment adjustment closure if differences in financing are taken into account.

\paragraph{The model.} The elastic-supply form of
Section~\ref{app:labour-closures}, \eqref{eq:closure-beta}, gives each sector
its own wage and supply curve and is solved in the unknowns
$X = (p, y, w(1{:}N), F) \in \mathbb{R}^{3N+1}$; the matrix's elastic-supply
row is this system at $\eta_s = 0.5$, and
Table~\ref{tab:sectoral} illustrates the ladder as well as its two ends.\footnote{The elasticity can differ
across sectors; writing $\eta_{s,i}$ for $\eta_s$ changes nothing in what
follows, but represents a venue for imposing further assumption on heterogenous rigidity or labor market conditions across sectors.} Because
the $N$ wages enter the $N$ zero-profit conditions through the price index
\eqref{eq:model-ip}, prices now depend on the $N$ employment levels and
Lemma~\ref{lem:price} no longer applies: demand reaches prices through the wage
structure, $p = \pi(A, w)$.

\begin{proposition}[The rigid corner is the $\eta = 0$ endpoint]\label{prop:nest}
Set $\eta_s = 0$. Then the elastic labor supply model and the $\eta = 0$
endpoint of the matrix have the same solution set.
\end{proposition}

\begin{proof}
Write both systems in the unknowns $X = (p, y, w, F)$ as four row blocks:
\begin{align*}
\text{(P)} \quad & p_i = c_i(p, w_i; A), & & i = 1,\dots,N,
   & & \text{\eqref{eq:zp}} \\
\text{(Y)} \quad & y_i = int_i + c_i + (1-m_i)(g^G_i + inv_i + g_i) + X_i,
   & & i = 1,\dots,N, & & \text{\eqref{eq:model-clearing}} \\
\text{(W)} \quad & \log L^{cm}_i(p,y_i,w_i) - \log \bar{L}_i
   - \eta_s \log \frac{w_i}{\mathcal{P}(p)} = 0, & & i = 1,\dots,N, & & \\
\text{(E)} \quad & \mathcal{P}(p) = 1. & & & &
\end{align*}
The quantity block (Y) draws labor income from the wage bill
$\sum_i w_i L_i$ inside \eqref{eq:cons} and from the tax rate $\tau$; this is
the only place, besides (W), where the two systems differ. The $\eta = 0$
endpoint is the same four blocks with two changes: (W) drops the supply term
($\eta_s = 0$), and (Y) evaluates the wage bill at the \emph{frozen}
allocation, $L_i := \bar{L}_i$, instead of at the cost-minimizing demand
$L^{cm}_i$.

\emph{(i)} Let $X$ solve the sectoral system at $\eta_s \equiv 0$. Its (W)
gives $L^{cm}_i(X) = \bar{L}_i$ for every $i$; substituting this into (Y) makes
that block equal the endpoint's (Y), while (P) and (E) are identical in the two
systems by inspection, so $X$ solves the endpoint system.
\emph{(ii)} Conversely, let $X$ solve the endpoint system. Its (W) gives
$L^{cm}_i(X) = \bar{L}_i$, which satisfies the sectoral (W) at
$\eta_s \equiv 0$, and substituting it into (Y) yields the sectoral (Y).
Both systems are square --- $3N+1$ unknowns, $3N+1$ equations --- and their
residual maps agree on the common solution set, so the solution sets coincide.
\end{proof}


\paragraph{The elastic end.} The results derived from the elastic labor supply model have a limit as $\eta_s\to\infty$,
and it does not directly correspond to the wage-pinned closure of the matrix. As indicated, while prices return to the baseline and the real wage to its anchor, the setups differ in the external transfer, which the elastic labor supply model determines and the matrix's row sets to zero. At the transfer the elastic labor supply model selects, the wage-pinned equations hold exactly; at the zero transfer they do not, and the difference is the transfer's own contribution to demand, as illustrated in table \ref{tab:sectoral} below. The limit therefore does not directly intersect with the results shown for the employment adjustment closure as $F$ remains unaligned across specifications

Pushing the uniform member past the executed cells, employment and consumption converge and the transfer turns positive; the table sets the
limit beside the closure at $F = 0$ and at $F = F^{\ast}$.

% Table B1 -- the sectoral family in three parts (rigid end / uniform ladder /
% elastic end).  Numbers from probe21_sectoral_family_limit.jl and probe23
% (2026-10-24), F2 financing, demand-only shocks, baseline elasticities
% theta = eps = 0.5, sigma = 0.9, eta = 1.
\begin{table}[htbp]
\centering
\caption{The sectoral family in three parts. \emph{Top:} the uniform ladder,
extended past the executed matrix to $\eta_s = 100$. \emph{Middle:} the rigid
end, where the family at $\eta_s = 0$ and the no-reallocation row coincide
to all printed digits. \emph{Bottom:} the elastic end. At the zero transfer the
family's limit and the employment-adjustment closure diverge; at the transfer the
family itself delivers, $F^{\mathrm{lim}}$, the two agree on the allocation ---
employment $1.002324$ against $1.002327$, consumption $-1.289$ against $-1.288$
--- and differ otherwise only in prices, which the closure sets exactly at the
baseline while the family's stand $8.6 \times 10^{-5}$ away. That residual is a finite
$\eta_s$ standing in for infinity, not a difference between the closures.
The columns are not in a common unit: prices, the largest absolute deviation of
any sectoral price from $1$; employment, the level relative to the no-program
baseline, so that $1.002327$ is $+0.233$ percent; consumption, the percentage
change from the baseline; and $F$, the net external transfer as a share of
baseline GDP. Financing F2, demand-only shocks.}
\label{tab:sectoral}
\small
\begin{tabular}{lrrrr}
\toprule
& $\max|p-1|$ & Employment & Consumption & $F$ \\
\midrule
\multicolumn{5}{@{}l}{\emph{Uniform elasticity}} \\
$\eta_s = 0.25$ & $0.152897$ & $1.000864$ & $-1.689$ & $-0.004424$ \\
$\eta_s = 0.5$  & $0.105666$ & $1.001255$ & $-1.571$ & $-0.002828$ \\
$\eta_s = 1$    & $0.065422$ & $1.001627$ & $-1.467$ & $-0.001444$ \\
$\eta_s = 2$    & $0.037171$ & $1.001914$ & $-1.391$ & $-0.000457$ \\
$\eta_s = 10$   & $0.008351$ & $1.002231$ & $-1.312$ & $+0.000562$ \\
$\eta_s = 100$  & $0.000859$ & $1.002318$ & $-1.291$ & $+0.000829$ \\
\midrule
\multicolumn{5}{@{}l}{\emph{The rigid end}} \\
Sectoral, $\eta_s = 0$    & $0.279009$ & $1.000000$ & $-1.980$ & $-0.008528$ \\
No reallocation ($\eta = 0$)  & $0.279009$ & $1.000000$ & $-1.980$ & $-0.008528$ \\
\midrule
\multicolumn{5}{@{}l}{\emph{The elastic end}} \\
Family limit ($\eta_s = 1000$)      & $0.000086$ & $1.002327$ & $-1.288$ & $+0.000856$ \\
Employment adjustment, $F = 0$      & $0.000000$ & $1.001260$ & $-1.533$ & $0.000000$ \\
Employment adjustment, $F = F^{\mathrm{lim}}$ & $0.000000$ & $1.002324$ & $-1.289$ & $+0.000856$ \\
\bottomrule
\end{tabular}
\end{table}


\emph{Two checks.} Solving the sectoral system at $\eta_s=0$ from the $\eta = 0$ solution reproduces all four blocks exactly, so no separate rigid
sectoral cell is reported: the corner \emph{is} the no-reallocation row of
Table~\ref{tab:results-matrix}. And nothing in the proof refers to $A\equiv\one$, so
the nesting survives supply shocks, under which the elastic end instead
separates and $\eta_s$ becomes identifiable --- whereas a demand-only
program leaves it a sensitivity parameter.
